% Geodesic Code Synthesis: Topology-Aware Program Generation
% via Manifold Navigation and Execution Prediction
%
% (C) 2026 Tristan Stoltz / Luminous Dynamics
% Draft - March 2026

\documentclass[11pt,a4paper]{article}

\usepackage[utf8]{inputenc}
\usepackage[T1]{fontenc}
\usepackage{amsmath,amssymb,amsthm}
\usepackage{algorithm}
\usepackage{algpseudocode}
\usepackage{graphicx}
\usepackage{hyperref}
\usepackage{booktabs}
\usepackage{xcolor}
\usepackage[margin=1in]{geometry}

\newtheorem{theorem}{Theorem}
\newtheorem{proposition}{Proposition}
\newtheorem{definition}{Definition}

\title{Geodesic Code Synthesis:\\
Topology-Aware Program Generation\\
via Manifold Navigation and Execution Prediction}

\author{Tristan Stoltz\\
Luminous Dynamics\\
\texttt{tristan.stoltz@evolvingresonantcocreationism.com}}

\date{March 2026 — Draft}

\begin{document}

\maketitle

\begin{abstract}
We introduce \textbf{Geodesic Code Synthesis (GCS)}, a fundamentally novel approach to
program generation that operates in the opposite direction from large language models.
Where LLMs generate code left-to-right via token prediction, GCS first determines the
\emph{topological structure} of the solution (number of loops, branches, recursive calls),
then composes algebraic combinators matching that topology, fills expression slots via
manifold interpolation, and verifies correctness through an $O(1)$ execution oracle
before emitting source code. GCS provides formal guarantees that LLMs fundamentally
cannot: topological invariants as hard constraints, pre-generation execution prediction
via Closed-form Continuous-time (CfC) neural dynamics, and sheaf-theoretic composability
verification. We implement GCS as a 4,849-line Rust crate built on Hyperdimensional
Computing (16,384-bit vectors) and demonstrate topology-correct skeleton generation
across 7 algorithmic patterns with 70 passing unit tests.
\end{abstract}

% ============================================================================
\section{Introduction}
% ============================================================================

Large language models generate code by predicting the next most likely token given
a prefix. This approach has achieved remarkable results on benchmarks like HumanEval
\cite{chen2021codex}, but suffers from fundamental limitations:

\begin{enumerate}
    \item \textbf{No structural awareness}: An LLM cannot guarantee that generated
    code contains exactly one loop, or that a recursive function has a base case.
    It generates tokens sequentially and hopes the structure works out.

    \item \textbf{No execution model}: LLMs have no mechanism to predict what code
    will \emph{do} before generating it. They cannot detect infinite loops, verify
    termination, or estimate computational complexity.

    \item \textbf{No composability guarantee}: Functions generated in isolation may
    have incompatible interfaces. LLMs have no formal mechanism to verify that
    locally-correct code assembles into a globally-correct program.

    \item \textbf{No structural equivalence}: LLMs treat \texttt{for i in 0..n}
    and \texttt{while i < n \{ i += 1 \}} as completely different token sequences,
    despite being structurally identical (both are $\beta_1 = 1$ loops).
\end{enumerate}

We propose \textbf{Geodesic Code Synthesis}, which addresses all four limitations
by generating programs \emph{structure-first}. GCS constructs the topological skeleton
of a program before any concrete expressions exist, ensuring structural correctness
by construction.

% ============================================================================
\section{Background}
% ============================================================================

\subsection{Hyperdimensional Computing}

Hyperdimensional Computing (HDC) \cite{kanerva2009hyperdimensional} represents
information as high-dimensional binary vectors (hypervectors). Three operations
form a complete algebra:
\begin{itemize}
    \item \textbf{Binding} ($\otimes$, XOR): Associates two concepts. Self-inverse:
    $A \otimes A = \mathbf{1}$.
    \item \textbf{Bundling} ($\oplus$, majority vote): Creates set representations.
    Approximately commutative.
    \item \textbf{Permutation} ($\rho^k$, cyclic shift): Encodes order/position.
\end{itemize}

We use 16,384-bit binary hypervectors ($\text{BinaryHV} \in \{-1, +1\}^{16384}$),
following the Symthaea architecture \cite{stoltz2026symthaea}.

\subsection{Closed-form Continuous-time Neural Networks}

CfC networks \cite{hasani2021liquid} provide $O(1)$ temporal dynamics via the
closed-form solution:
\begin{equation}
    x(t + \Delta t) = x_\infty + (x(t) - x_\infty) \cdot \exp(-\Delta t / \tau)
\end{equation}
where $x_\infty$ is the equilibrium (attractor) state and $\tau$ is the time constant.
This allows predicting the state after $N$ steps without simulating each step —
the key enabler of the execution oracle.

\subsection{Persistent Homology}

Persistent homology \cite{edelsbrunner2002topological,carlsson2009topology} computes
topological features of data across scales. The Betti numbers $\beta_k$ count
$k$-dimensional ``holes'':
\begin{itemize}
    \item $\beta_0$: connected components
    \item $\beta_1$: independent cycles (loops)
    \item $\beta_2$: enclosed voids
\end{itemize}

Applied to a Program Dependence Graph, these capture structural properties:
$\beta_0$ measures module coupling, $\beta_1$ counts loops and recursive cycles,
and $\beta_2$ detects unreachable code regions.

\subsection{Sheaf Theory}

A sheaf \cite{kashiwara1990sheaves} formalizes ``local-to-global'' consistency:
data defined on overlapping open sets must agree on their intersections.
Applied to code, local sections are function implementations, open sets are
connected subgraphs of the call graph, and the gluing condition requires
that caller expectations match callee interfaces.

% ============================================================================
\section{Architecture}
% ============================================================================

GCS consists of six layers, each constraining the next:

\begin{enumerate}
    \item \textbf{Program Dependence Graph (PDG)}: Merges control-flow and data
    dependency information from parsed ASTs into a single directed graph.

    \item \textbf{Topological Fingerprint}: Computes Betti numbers
    $(\beta_0, \beta_1, \beta_2)$ from the PDG's simplicial complex (control-flow
    edges only, excluding data dependencies to prevent spurious cycles).

    \item \textbf{Program Manifold}: A fiber bundle where the base space is
    (algorithm pattern $\times$ type signature $\times$ topological fingerprint)
    and each fiber contains concrete implementations sharing the same structural
    properties.

    \item \textbf{Execution Oracle}: CfC neurons simulate program execution with
    type-aware time constants ($\tau_{\text{arithmetic}} = 0.1$,
    $\tau_{\text{loop}} = 1.0$, $\tau_{\text{IO}} = 10.0$). Predicts outcomes
    and termination in $O(1)$ via the closed-form temporal jump.

    \item \textbf{Sheaf Coherence}: Verifies that locally-correct functions
    compose into a globally-correct program via HDC similarity between
    caller-expected and callee-actual interfaces.

    \item \textbf{Geodesic Synthesis}: Walks the shortest path through the
    program manifold from specification to implementation, constrained by
    topology, guided by execution prediction, and verified by sheaf coherence.
\end{enumerate}

\subsection{Topological Skeleton Synthesis}

The key innovation is \emph{skeleton-first generation}. Each algebraic combinator
has a known topological signature:

\begin{table}[h]
\centering
\begin{tabular}{lcc}
\toprule
\textbf{Combinator} & $\Delta\beta_1$ & \textbf{Code Pattern} \\
\midrule
\texttt{Sequence} & 0 & Linear execution \\
\texttt{Branch} & 0 & \texttt{if/else}, \texttt{match} \\
\texttt{Iterate} & +1 & \texttt{for}, \texttt{while}, \texttt{loop} \\
\texttt{Recurse} & +1 & Recursive function call \\
\texttt{Map/Filter/Reduce} & +1 & Iterator combinators \\
\texttt{Leaf} & 0 & Expression (no structure) \\
\bottomrule
\end{tabular}
\caption{Topological signatures of skeleton combinators. Composition is additive:
the topology of a composed skeleton is the sum of its parts.}
\label{tab:signatures}
\end{table}

\begin{proposition}[Topological Correctness by Construction]
If a skeleton $S$ is composed from combinators $c_1, \ldots, c_n$ and the target
topology requires $\beta_1 = k$ cycles, then
$\sum_{i=1}^{n} \Delta\beta_1(c_i) = k$ guarantees the skeleton has exactly $k$
loops, regardless of how the expression slots are filled.
\end{proposition}

\begin{proof}
Each combinator adds a fixed number of back-edges to the control-flow graph.
The simplicial complex constructed from control-flow edges only has
$\beta_1 = |E_{\text{back}}| = \sum \Delta\beta_1(c_i)$ by the Euler characteristic
relation $\beta_1 = |E| - |V| + \beta_0$ applied to the 1-skeleton, where
back-edges are the only source of cycles in a DAG with added back-edges.
\end{proof}

\subsection{$O(1)$ Execution Prediction}

The execution oracle repurposes CfC dynamics for program state evolution.
Given a program state $s(t)$ and a statement encoded as hypervector $h$:

\begin{equation}
    s_\infty = s(t) \otimes h \quad \text{(HDC binding = state transformation)}
\end{equation}
\begin{equation}
    s(t+1) = s_\infty + (s(t) - s_\infty) \cdot \exp(-1/\tau_{\text{op}})
\end{equation}

For a loop with $N$ iterations, the $O(1)$ prediction:
\begin{equation}
    s(t+N) = s_\infty + (s(t) - s_\infty) \cdot \exp(-N/\tau)
\end{equation}

\textbf{Convergence criterion}: A loop terminates if
$\|s(t+N) - s_\infty\| < \epsilon$, which occurs when $N > \tau \cdot \ln(1/\epsilon)$.
The time constant $\tau$ thus gives an approximate complexity measure:
small $\tau$ implies fast convergence (constant/logarithmic time),
large $\tau$ implies slow convergence (linear/quadratic).

% ============================================================================
\section{Implementation}
% ============================================================================

GCS is implemented as a Rust crate (\texttt{symthaea-geodesic}) within the
Symthaea workspace. Table~\ref{tab:modules} summarizes the implementation.

\begin{table}[h]
\centering
\begin{tabular}{llrr}
\toprule
\textbf{Layer} & \textbf{Module} & \textbf{LOC} & \textbf{Tests} \\
\midrule
1. PDG & \texttt{pdg.rs} & 742 & 9 \\
2. Topology & \texttt{topology.rs} & 510 & 7 \\
3. Manifold & \texttt{manifold.rs} & 374 & 6 \\
4. Oracle & \texttt{execution\_oracle.rs} & 688 & 9 \\
5. Sheaf & \texttt{sheaf.rs} & 476 & 9 \\
6. Synthesis & \texttt{synthesis.rs} & 731 & 9 \\
7. Skeleton & \texttt{skeleton\_synthesis.rs} & 857 & 12 \\
— & Bridge/bootstrap & 419 & 9 \\
\midrule
\textbf{Total} & & \textbf{4,849} & \textbf{70} \\
\bottomrule
\end{tabular}
\caption{GCS implementation summary. All 70 tests pass.}
\label{tab:modules}
\end{table}

% ============================================================================
\section{Evaluation}
% ============================================================================

\subsection{Skeleton Topology Correctness}

We test skeleton generation across 7 algorithmic patterns, each with a specified
target topology. Results in Table~\ref{tab:skeleton_results}.

\begin{table}[h]
\centering
\begin{tabular}{lccc}
\toprule
\textbf{Pattern} & \textbf{Target $\beta_1$} & \textbf{Generated $\beta_1$} & \textbf{Match} \\
\midrule
Simple addition & 0 & 0 & \checkmark \\
Bubble sort (nested) & 2 & 2 & \checkmark \\
Binary search & 1 & 1 & \checkmark \\
Fibonacci (recursive) & 1 & 1 & \checkmark \\
Filter (iterator) & 1 & 1 & \checkmark \\
Reduce (fold) & 1 & 1 & \checkmark \\
Branch only & 0 & 0 & \checkmark \\
\midrule
\textbf{Accuracy} & & & \textbf{7/7 (100\%)} \\
\bottomrule
\end{tabular}
\caption{Skeleton topology correctness. All generated skeletons match target Betti numbers.}
\label{tab:skeleton_results}
\end{table}

\subsection{PDG Analysis of Real Code}

We analyze a binary search implementation (16 lines of Rust) via the PDG pipeline:

\begin{table}[h]
\centering
\begin{tabular}{lccc}
\toprule
\textbf{Metric} & \textbf{Before Calibration} & \textbf{After Calibration} & \textbf{Expected} \\
\midrule
Loop count & 0 & 1 & 1 \\
$\beta_1$ & 9 & 2 & 1--2 \\
$\beta_2$ & 8 & 0 & 0 \\
\bottomrule
\end{tabular}
\caption{PDG topology calibration. Filtering to control-flow edges eliminates
spurious cycles from data dependency edges.}
\label{tab:calibration}
\end{table}

The initial analysis ($\beta_1 = 9$) was caused by including data dependency (def-use)
edges in the simplicial complex, which created spurious 1-cycles. After restricting to
control-flow edges only, the topology correctly identifies one loop ($\beta_1 = 2$
includes the loop back-edge plus one branch-merge cycle within the loop body).

% ============================================================================
\section{Capabilities Beyond LLMs}
% ============================================================================

Table~\ref{tab:capabilities} summarizes what GCS enables that token-prediction
models fundamentally cannot provide.

\begin{table}[h]
\centering
\begin{tabular}{lcc}
\toprule
\textbf{Capability} & \textbf{LLM} & \textbf{GCS} \\
\midrule
Pre-generation execution prediction & \texttimes & \checkmark \\
Topological invariants as hard constraints & \texttimes & \checkmark \\
$O(1)$ loop termination prediction & \texttimes & \checkmark \\
Formal composability (sheaf coherence) & \texttimes & \checkmark \\
Structural equivalence across syntax & \texttimes & \checkmark \\
Complexity estimation from dynamics & \texttimes & \checkmark \\
Causal intervention (``what if?'') & \texttimes & \checkmark \\
Open-ended generative synthesis & \checkmark & \texttimes \\
Natural language fluency & \checkmark & \texttimes \\
Memorized code corpus & \checkmark & \texttimes \\
\bottomrule
\end{tabular}
\caption{Capability comparison. GCS and LLMs have complementary strengths;
the optimal system uses GCS for verification and structural synthesis,
with LLMs for expression-level generation.}
\label{tab:capabilities}
\end{table}

% ============================================================================
\section{Discussion}
% ============================================================================

GCS is not a replacement for LLMs — it is a \emph{complementary architecture}
that provides formal guarantees LLMs cannot. The optimal coding system uses
GCS for three purposes:

\begin{enumerate}
    \item \textbf{Verification gate}: After LLM generation, GCS builds a PDG,
    checks topology, runs the oracle, and verifies sheaf coherence. Violations
    are fed back as hard constraints for regeneration.

    \item \textbf{Structure-first generation}: For problems where the topology
    is known (``implement binary search'' $\Rightarrow$ $\beta_1 = 1$), GCS
    generates the correct skeleton before the LLM fills in expressions.

    \item \textbf{Manifold-guided retrieval}: Instead of generating from scratch,
    query the program manifold for the nearest fiber to the specification and
    adapt existing implementations.
\end{enumerate}

\subsection{Limitations}

The execution oracle approximates program execution — it cannot precisely predict
outputs for complex programs. The PDG construction from source code is heuristic
(line-level scanning rather than full AST analysis). The manifold requires
population from an existing codebase to be useful.

These limitations are addressable: tree-sitter integration for production PDGs,
oracle calibration against ground truth, and automatic codebase indexing.

% ============================================================================
\section{Related Work}
% ============================================================================

\textbf{Program synthesis}: Gulwani et al. \cite{gulwani2017program} survey
synthesis from specifications. GCS differs by using topological constraints
rather than logical specifications.

\textbf{Topological data analysis for software}: Persistence-based analysis
of code has been explored for bug detection \cite{dlotko2019topological},
but not for code generation.

\textbf{Neural program synthesis}: Devlin et al. \cite{devlin2017robustfill}
and Parisotto et al. \cite{parisotto2017neuro} use neural networks for program
synthesis but without topological structure.

\textbf{Algebra of Programming}: Bird and de Moor \cite{bird1997algebra}
formalize program transformation as algebraic operations. GCS extends this
with HDC encoding and topological constraints.

% ============================================================================
\section{Conclusion}
% ============================================================================

Geodesic Code Synthesis demonstrates that program generation can be approached
as navigation through a topological manifold rather than token prediction.
By determining structure first (Betti numbers), composing algebraic skeletons
(combinators with known topological signatures), predicting execution ($O(1)$
CfC dynamics), and verifying composability (sheaf coherence), GCS provides
formal guarantees that complement LLM-based generation.

The 4,849-line implementation with 70 passing tests validates the approach.
Skeleton generation achieves 100\% topology correctness across 7 algorithmic
patterns, and the PDG analysis pipeline produces meaningful structural
characterizations of real Rust code after calibration.

\bibliographystyle{plain}
\begin{thebibliography}{99}

\bibitem{kanerva2009hyperdimensional}
P.~Kanerva.
\newblock Hyperdimensional computing: An introduction to computing in
  distributed representation with high-dimensional random vectors.
\newblock {\em Cognitive Computation}, 1(2):139--159, 2009.

\bibitem{hasani2021liquid}
R.~Hasani, M.~Lechner, A.~Amini, D.~Rus, and R.~Grosu.
\newblock Liquid time-constant networks.
\newblock In {\em AAAI}, 2021.

\bibitem{edelsbrunner2002topological}
H.~Edelsbrunner, D.~Letscher, and A.~Zomorodian.
\newblock Topological persistence and simplification.
\newblock {\em Discrete \& Computational Geometry}, 28:511--533, 2002.

\bibitem{carlsson2009topology}
G.~Carlsson.
\newblock Topology and data.
\newblock {\em Bulletin of the AMS}, 46(2):255--308, 2009.

\bibitem{kashiwara1990sheaves}
M.~Kashiwara and P.~Schapira.
\newblock {\em Sheaves on Manifolds}.
\newblock Springer, 1990.

\bibitem{chen2021codex}
M.~Chen et~al.
\newblock Evaluating large language models trained on code.
\newblock {\em arXiv:2107.03374}, 2021.

\bibitem{stoltz2026symthaea}
T.~Stoltz.
\newblock Symthaea: A holographic liquid brain for consciousness-first AI.
\newblock Technical report, Luminous Dynamics, 2026.

\bibitem{gulwani2017program}
S.~Gulwani, O.~Polozov, and R.~Singh.
\newblock Program synthesis.
\newblock {\em Foundations and Trends in Programming Languages}, 4(1--2):1--119, 2017.

\bibitem{dlotko2019topological}
P.~D{\l}otko, R.~Ghrist, and Y.~Kang.
\newblock Topological approaches to analysis of software.
\newblock {\em Applied Algebraic Topology Network}, 2019.

\bibitem{devlin2017robustfill}
J.~Devlin, J.~Uesato, S.~Bhupatiraju, R.~Singh, A.~Mohamed, and P.~Kohli.
\newblock RobustFill: Neural program learning under noisy I/O.
\newblock In {\em ICML}, 2017.

\bibitem{parisotto2017neuro}
E.~Parisotto, A.~Mohamed, R.~Singh, L.~Li, D.~Zhou, and P.~Kohli.
\newblock Neuro-symbolic program synthesis.
\newblock In {\em ICLR}, 2017.

\bibitem{bird1997algebra}
R.~Bird and O.~de~Moor.
\newblock {\em Algebra of Programming}.
\newblock Prentice Hall, 1997.

\end{thebibliography}

\end{document}
